\documentclass[preview]{standalone}

\usepackage{amsmath}
\usepackage{amssymb}
\usepackage{stellar}
\usepackage{definitions}
\usepackage{bettelini}

\begin{document}

\id{analytic-functions}
\genpage

\section{Taylor series and analyticity}

\begin{snippettheorem}{borel-theorem-smooth-jets}{Borel's theorem}
    Let \(\{b_n\}_{n\geq0}\) be any real sequence, let \((a,b)\) be an open
    interval, and let \(x_0\in(a,b)\). There exists
    \(f\in\mathcal C^\infty((a,b))\) such that
    \[
        f^{(n)}(x_0)=b_n,
        \qquad \forall n\geq0.
    \]
\end{snippettheorem}

\begin{snippetexample}{borel-theorem-zero-radius-taylor-series-example}{A smooth function with a Taylor series of zero radius}
    In Borel's theorem, choose \(b_n=(n!)^2\). The resulting smooth function
    has Taylor coefficients
    \[
        \frac{b_n}{n!}=n!,
    \]
    so its Taylor series at \(x_0\) has radius of convergence \(0\). Thus
    smoothness alone imposes no positive radius of convergence on a Taylor
    series.
\end{snippetexample}

\begin{snippetdefinition}{real-analytic-function-definition}{Real analytic function}
    Let \((a,b)\subseteq\realnumbers\). A \function
    \(f\in\mathcal C^\infty((a,b))\) is \emph{real analytic} if, for every
    \(x_0\in(a,b)\), there exists a neighborhood \(U\subseteq(a,b)\) of
    \(x_0\) such that
    \[
        f(x)=\sum_{n=0}^{\infty}\frac{f^{(n)}(x_0)}{n!}(x-x_0)^n,
        \qquad \forall x\in U.
    \]
    The class of real analytic functions on \((a,b)\) is denoted by
    \(\mathcal C^\omega((a,b))\).
\end{snippetdefinition}

\begin{snippet}{real-analytic-functions-local-character}
    Analyticity is local: the Taylor series at each point only needs to agree
    with the function on some neighborhood of that point, not necessarily on
    the entire domain. Local power-series representations agree on overlaps.
    Polynomials are the simplest analytic functions.
\end{snippet}

\begin{snippetexample}{geometric-function-taylor-radius-by-center-example}{Taylor radius depends on the center}
    The \function
    \[
        f(x)=\frac1{1-x}
    \]
    is real analytic on \(( -\infty,1)\), although its Maclaurin expansion
    \[
        \frac1{1-x}=\sum_{n=0}^{\infty}x^n
    \]
    has radius \(1\). Centered at \(x_0=1/3\), its Taylor series has radius
    \(2/3\), the distance from the center to the singularity at \(1\).
\end{snippetexample}

\begin{snippetexample}{rational-function-complex-singularities-taylor-radius-example}{Complex singularities determine a real Taylor radius}
    The real \function
    \[
        f(x)=\frac1{1+x^2}
    \]
    is real analytic on all of \(\realnumbers\), but its Maclaurin series
    \[
        \sum_{n=0}^{\infty}(-1)^nx^{2n}
    \]
    has radius \(1\). In the complex plane, the nearest singularities are the
    poles at \(\pm i\). If the series is centered at \(x_0=2\), its radius is
    \(\sqrt5\), the distance from \(2\) to either pole.
\end{snippetexample}

\section{Sufficient derivative bounds}

\begin{snippettheorem}{analyticity-from-exponential-derivative-bound-theorem}{Analyticity from an exponential derivative bound}
    Let \(f\in\mathcal C^\infty((a,b))\). Suppose there exist constants
    \(M,L>0\) such that
    \[
        |f^{(n)}(x)|\leq ML^n,
        \qquad \forall x\in(a,b),\ \forall n\geq0.
    \]
    Then \(f\) is analytic on \((a,b)\). More precisely, for every
    \(x_0,x\in(a,b)\),
    \[
        f(x)=\sum_{n=0}^{\infty}\frac{f^{(n)}(x_0)}{n!}(x-x_0)^n.
    \]
\end{snippettheorem}

\begin{snippetproof}{analyticity-from-exponential-derivative-bound-theorem-proof}{analyticity-from-exponential-derivative-bound-theorem}{Analyticity from an exponential derivative bound}
    Fix \(x_0,x\in(a,b)\). Taylor's formula with Lagrange remainder gives,
    for every \(k\), a point \(c_k\) between \(x_0\) and \(x\) such that
    \[
        f(x)=\sum_{n=0}^{k}\frac{f^{(n)}(x_0)}{n!}(x-x_0)^n
        +\frac{f^{(k+1)}(c_k)}{(k+1)!}(x-x_0)^{k+1}.
    \]
    The remainder satisfies
    \[
        \left|\frac{f^{(k+1)}(c_k)}{(k+1)!}(x-x_0)^{k+1}\right|
        \leq M\frac{(L|x-x_0|)^{k+1}}{(k+1)!}\longrightarrow0.
    \]
    Letting \(k\to\infty\) proves the representation.
\end{snippetproof}

\begin{snippet}{analyticity-from-exponential-derivative-bound-local-use}
    The derivative bound is sufficient, not necessary, and it may be applied
    locally. For example, \(1/(1-x)\) is analytic on \(( -\infty,1)\), but no
    single pair \(M,L\) satisfies this bound on the whole interval.
\end{snippet}

\begin{snippetproposition}{equibounded-derivatives-imply-analyticity}{Equibounded derivatives imply analyticity}
    Let \(f\in\mathcal C^\infty((a,b))\). If there exists \(K>0\) such that
    \[
        |f^{(n)}(x)|\leq K,
        \qquad \forall x\in(a,b),\ \forall n\geq0,
    \]
    then \(f\) is analytic on \((a,b)\).
\end{snippetproposition}

\begin{snippetproof}{equibounded-derivatives-imply-analyticity-proof}{equibounded-derivatives-imply-analyticity}{Equibounded derivatives imply analyticity}
    Apply the preceding theorem with \(M=K\) and \(L=1\).
\end{snippetproof}

\begin{snippetexample}{elementary-functions-analyticity-from-derivative-bounds}{Analyticity of elementary functions}
    The derivatives of sine and cosine are uniformly bounded on
    \(\realnumbers\), so both functions are real analytic. The derivatives of
    \(e^x\) equal \(e^x\); on every interval \(( -\infty,A)\) they are
    bounded by \(e^A\). Hence the exponential is real analytic on
    \(\realnumbers\), and its power series has infinite radius, defining an
    entire function on \(\complexnumbers\).
\end{snippetexample}

\section{Binomial series}

\includesnpt{binomial-maclaurin-series}

\begin{snippet}{generalized-binomial-series-radius}
    If \(\alpha\in\naturalnumbers\cup\{0\}\), the series terminates. For
    every other \(\alpha\), its radius of convergence is \(1\).
\end{snippet}

\begin{snippetproof}{binomial-maclaurin-series-proof}{binomial-maclaurin-series}{Binomial Maclaurin series}
    If \(\alpha\) is a nonnegative integer, the identity is the ordinary
    binomial theorem. Otherwise, the coefficient ratio satisfies
    \[
        \left|\frac{\binom{\alpha}{n}}
        {\binom{\alpha}{n+1}}\right|
        =\frac{n+1}{|\alpha-n|}\longrightarrow1,
    \]
    so the radius is \(1\). Let
    \[
        g(x)=\sum_{n=0}^{\infty}\binom{\alpha}{n}x^n.
    \]
    Termwise differentiation and the recurrence
    \((n+1)\binom{\alpha}{n+1}=(\alpha-n)\binom{\alpha}{n}\) give
    \[
        (1+x)g'(x)=\alpha g(x),
        \qquad g(0)=1.
    \]
    On \((-1,1)\), the unique solution of this initial-value problem is
    \(g(x)=(1+x)^\alpha\), proving the identity.
\end{snippetproof}

\section{Power-series approximation}

\begin{snippet}{taylor-versus-stone-weierstrass-approximation}
    Stone--Weierstrass says that continuous functions on compact subsets of
    the real line can be uniformly approximated by polynomials. Taylor
    approximation is more rigid: for an analytic function, improving the
    approximation amounts to adding further terms whose coefficients are
    already fixed by derivatives at the center. A general
    Stone--Weierstrass approximation has no such nested structure and may need
    to be rebuilt at each degree.
\end{snippet}

\end{document}
